\chapter{Umkehrfunktionen und Komposition}
\label{b08guide:ch:inverses}

\section{Rechenregeln für Umkehrfunktionen}

\begin{theorem}[Beidseitige Umkehrgesetze]
\label{b08guide:thm:inverselaws}
Ist $f\colon A\bij B$, so gelten
\[
 \begin{gathered}
 f^{-1}\colon B\bij A,\qquad (f^{-1})^{-1}=f,\\
 f^{-1}\circ f=\operatorname{id}_A,
 \qquad f\circ f^{-1}=\operatorname{id}_B.
 \end{gathered}
\]
Für jede Menge $A$ ist $\operatorname{id}_A$ bijektiv, und
$(\operatorname{id}_A)^{-1}=\operatorname{id}_A$.
\end{theorem}

\begin{theorem}[Komposition und Umkehrung der Reihenfolge]
\label{b08guide:thm:composition}
Für $f\colon A\bij B$ und $g\colon B\bij C$ gilt
\[
 g\circ f\colon A\bij C,
 \qquad (g\circ f)^{-1}=f^{-1}\circ g^{-1}.
\]
Die inverse Komposition wird somit auf $C$ zuerst mit $g^{-1}$
und danach mit $f^{-1}$ ausgewertet.
\end{theorem}

\begin{theorem}[Zwei-aus-drei-Regel]
\label{b08guide:thm:twoofthree}
Seien $f\colon A\to B$ und $g\colon B\to C$ Funktionen.
Sind zwei der drei Funktionen
\[
 f,\qquad g,\qquad g\circ f
\]
bijektiv zwischen ihren angegebenen Mengen, so ist es auch die dritte.
Insbesondere gelten
\[
 \begin{aligned}
 f\colon A\bij B&\ \Longrightarrow\
 (g\circ f\text{ bijektiv}\ \Longleftrightarrow\ g\text{ bijektiv}),\\
 g\colon B\bij C&\ \Longrightarrow\
 (g\circ f\text{ bijektiv}\ \Longleftrightarrow\ f\text{ bijektiv}).
 \end{aligned}
\]
\end{theorem}

\section{Eindeutiges Lösen von Funktionsgleichungen}

\begin{theorem}[Gleichungen mit einer gegebenen Bijektion]
\label{b08guide:thm:equations}
Sei $f\colon A\bij B$ und $T$ eine beliebige Menge.
\begin{enumerate}
\item Zu jeder Funktion $k\colon T\to B$ gibt es genau eine Funktion
      $h\colon T\to A$ mit $f\circ h=k$. Sie lautet $h=f^{-1}\circ k$.
\item Zu jeder Funktion $\ell\colon A\to T$ gibt es genau eine Funktion
      $q\colon B\to T$ mit $q\circ f=\ell$. Sie lautet $q=\ell\circ f^{-1}$.
\end{enumerate}
Allgemeiner seien $u\colon A\bij B$, $v\colon C\bij D$ und
$k\colon A\to D$ gegeben. Dann besitzt
\[
 v\circ h\circ u=k
\]
genau eine Lösung $h\colon B\to C$, nämlich
\[
 h=v^{-1}\circ k\circ u^{-1}.
\]
Diese Lösung ist genau dann bijektiv, wenn $k$ bijektiv ist.
\end{theorem}

\section{Selbstbijektionen ohne zusätzliche Struktur}

\begin{definition}[Permutation einer Menge]
Eine \emph{Permutation} von $A$ ist eine Bijektion $p\colon A\bij A$.
Eine \emph{Involution} ist eine Funktion $p\colon A\to A$ mit
$p\circ p=\operatorname{id}_A$; die Identität ist dabei zugelassen.
\end{definition}

\begin{theorem}[Selbstumkehrung]
\label{b08guide:thm:selfinverse}
Für eine Funktion $p\colon A\to A$ sind äquivalent:
\[
 p\circ p=\operatorname{id}_A
 \quad\Longleftrightarrow\quad
 p\text{ ist bijektiv und }p^{-1}=p.
\]
\end{theorem}

\section{Einen Funktionswert vorschreiben}

\begin{theorem}[Anpassung einer vorhandenen Bijektion]
\label{b08guide:thm:pointed}
Seien $f\colon A\bij B$, $a\in A$ und $b\in B$ gegeben.
Dann gibt es eine Bijektion $g\colon A\bij B$ mit $g(a)=b$, die
außerhalb von höchstens zwei Argumenten mit $f$ übereinstimmt.
Ist $f(a)=b$, kann $g=f$ gewählt werden. Andernfalls ist
$c=f^{-1}(b)\ne a$, und eine solche Bijektion ist
\[
 g(x)=\begin{cases}
 b,&x=a,\\
 f(a),&x=c,\\
 f(x),&x\in A\setminus\{a,c\}.
 \end{cases}
\]
\end{theorem}
